In this appendix we explain why part~(3) of
\Cref{theorem:cyclic_cases_after_blowup} is difficult to prove by induction
on $r$ from the degree-$p$ case, part~(2) of the same theorem.



\section{Quotient Towers of Symmetric-Power Torsors}\label{subsection:quotient_torsors_towers}

We now study how the symmetric power construction of \Cref{prop:symmetric_power_torsor} interacts with the canonical decomposition of \Cref{prop:quotient_torsors_finite}.
Throughout this section, $G$ denotes a finite abelian group, $H \subseteq G$ a subgroup, and $\cP \to X$ a $G$-torsor over an $S$-scheme $X$ that is Zariski-locally quasi-projective over $S$.
Recall that $\cP \to X$ factors canonically as
\[
\cP \xrightarrow{\phi} \cP/H \xrightarrow{\psi} X,
\]
where $\cP/H = (G/H) \times^G \cP$ is a $G/H$-torsor over $X$ and $\phi$ is an $H$-torsor.

In addition to the kernel $\KG = \ker(G^d \to G)$ of the summation morphism appearing in \Cref{prop:symmetric_power_torsor}, we shall need the analogous \emph{summation kernels}
\[
K_H := \ker(H^d \to H), \qquad K_{G/H} := \ker((G/H)^d \to G/H).
\]
Applying the snake lemma to the commutative diagram
\[
\begin{tikzcd}[row sep=large, column sep=large]
0 \arrow[r] & H^d \arrow[r] \arrow[d, twoheadrightarrow] & G^d \arrow[r] \arrow[d, twoheadrightarrow] & (G/H)^d \arrow[r] \arrow[d, twoheadrightarrow] & 0 \\
0 \arrow[r] & H \arrow[r] & G \arrow[r] & G/H \arrow[r] & 0,
\end{tikzcd}
\]
in which the vertical maps are the surjective summation morphisms, yields a short exact sequence
\begin{equation}\label{eq:summation_kernels_ses}
0 \to K_H \to \KG \to K_{G/H} \to 0.
\end{equation}
The induced isomorphism $\KG/K_H \cong K_{G/H}$ will play a central role below.

We define four quotient schemes
\begin{align*}
    \cP^{[d]_G} &:= (\KG \rtimes S_d) \backslash \cP^{\times_S d}, \\
    \cP^{[d]_H} &:= (K_H \rtimes S_d) \backslash \cP^{\times_S d}, \\
    (\cP/H)^{(d)} &:= (H^d \rtimes S_d) \backslash \cP^{\times_S d}, \\
    (\cP/H)^{[d]} &:= ((H^d + \KG) \rtimes S_d) \backslash \cP^{\times_S d},
\end{align*}
where $S_d$ acts on $\cP^{\times_S d}$ as in \Cref{cor:sym-power-quotient} and the subgroups $K_H \subseteq \KG \subseteq H^d + \KG$ and $H^d \subseteq H^d + \KG$ of $G^d$ act diagonally on $\cP^{\times_S d}$ via the $G$-action on $\cP$. The first space is the symmetric $G$-power of $\cP$ as in \Cref{prop:symmetric_power_torsor}; the remaining three are, respectively, the symmetric $H$-power of $\cP$ relative to $\phi$, the symmetric power of the scheme $\cP/H$, and the symmetric $G/H$-power of $\cP/H$.
\textcolor{red}{\textbf{The latter two identifications follow from $(\cP/H)^{\times_S d} \cong H^d \backslash \cP^{\times_S d}$ together with the fact that the preimage of $K_{G/H} \subseteq (G/H)^d$ under $G^d \twoheadrightarrow (G/H)^d$ equals $H^d + \KG$.}}
\textcolor{red}{the last isomorphism like that, is from the correspondence theorem, claude generated at the end of the article, add it somewhere}

\begin{prop}\label{prop:torsor_tower_diagram}
The four quotient schemes above exist and fit into a canonical commutative diagram
\[
\begin{tikzcd}[row sep=large, column sep=large]
    \cP^{[d]_H} \arrow[r, "H"] \arrow[d] & (\cP/H)^{(d)} \arrow[d, "q"] \\
    \cP^{[d]_G} \arrow[r, "H"'] \arrow[dr, "G"', bend right=20] & (\cP/H)^{[d]} \arrow[d, "G/H"] \\
    & X^{(d)},
\end{tikzcd}
\]
in which each labelled arrow is a torsor under the indicated constant group. The right-hand vertical $q$ is the canonical projection of \Cref{cor:sym-power-quotient}, applied to the $G/H$-torsor $\cP/H \to X$; in general it is \emph{not} a torsor under any constant group scheme (see the remark below). The unlabelled left-hand vertical $\cP^{[d]_H} \to \cP^{[d]_G}$ is the canonical projection induced by the inclusion of $S_d$-stable subgroups $K_H \rtimes S_d \subseteq \KG \rtimes S_d$ in $G^d \rtimes S_d$; it is the base change of $q$ along the structural map $\cP^{[d]_G} \to (\cP/H)^{[d]}$ via the Cartesian property below, and likewise fails to be a torsor under a constant group scheme in general.

Moreover, the upper square is Cartesian:
\[
\cP^{[d]_H} \cong \cP^{[d]_G} \times_{(\cP/H)^{[d]}} (\cP/H)^{(d)}.
\]
\end{prop}

\begin{proof}
The four labelled arrows are torsors under the indicated constant groups by direct application of \Cref{prop:symmetric_power_torsor} and \Cref{prop:quotient_torsors_finite}:
\begin{itemize}
    \item Applied to the $G$-torsor $\cP \to X$, \Cref{prop:symmetric_power_torsor} gives the $G$-torsor $\cP^{[d]_G} \to X^{(d)}$ (the diagonal arrow).
    \item Applied to the $H$-torsor $\phi: \cP \to \cP/H$ (taking $\cP/H$ in place of the base $X$), \Cref{prop:symmetric_power_torsor} gives the $H$-torsor $\cP^{[d]_H} \to (\cP/H)^{(d)}$ (the top arrow).
    \item Applied to the $G/H$-torsor $\psi: \cP/H \to X$, \Cref{prop:symmetric_power_torsor} gives the $G/H$-torsor $(\cP/H)^{[d]} \to X^{(d)}$ (the bottom-right vertical), and \Cref{cor:sym-power-quotient} produces the canonical morphism $q : (\cP/H)^{(d)} \to (\cP/H)^{[d]}$.
\end{itemize}

The factorization $\cP^{[d]_G} \xrightarrow{H} (\cP/H)^{[d]} \xrightarrow{G/H} X^{(d)}$ of the diagonal arrow into an $H$-torsor followed by a $G/H$-torsor is the canonical decomposition of \Cref{prop:quotient_torsors_finite}, applied to the $G$-torsor $\cP^{[d]_G} \to X^{(d)}$ and the subgroup $H \subseteq G$.
Indeed, both $\cP^{[d]_G} \times^G (G/H)$ and $(\cP/H)^{[d]}$ are realized as the admissible quotient
$((H^d + \KG) \rtimes S_d) \backslash \cP^{\times_S d}$.
\textcolor{red}{here maybe add more info, like isomorphism theorem for torsors and qoutients, or something like that, to make it more clear, from which it follows}

To prove the upper square is Cartesian, set $F := \cP^{[d]_G} \times_{(\cP/H)^{[d]}} (\cP/H)^{(d)}$.
Base-changing the $H$-torsor $\cP^{[d]_G} \to (\cP/H)^{[d]}$ along $q$ exhibits $F$ as an $H$-torsor over $(\cP/H)^{(d)}$.
The universal property of the fiber product furnishes a canonical morphism $\cP^{[d]_H} \to F$ over $(\cP/H)^{(d)}$,
which is $H$-equivariant since the $H$-actions on both source and target are induced by the diagonal $H$-action on $\cP^{\times_S d}$ \textcolor{red}{add earlier in the section of torsors that in fact base changing an $H$-torsor along a morphism gives another $H$-torsor
 is a torsor}.
 Any morphism of $H$-torsors is an isomorphism, so $\cP^{[d]_H} \cong F$ and the upper square is Cartesian. The left vertical $\cP^{[d]_H} \to \cP^{[d]_G}$ is then identified, via this Cartesian square, with the base change of $q$ along $\cP^{[d]_G} \to (\cP/H)^{[d]}$.
\end{proof}

\begin{rem*}
The two vertical morphisms $q : (\cP/H)^{(d)} \to (\cP/H)^{[d]}$ and $\cP^{[d]_H} \to \cP^{[d]_G}$ are not, in general, torsors under any constant group scheme. The obstruction is purely group-theoretic: in $\KG \rtimes S_d$ the subgroup $K_H \rtimes S_d$ is not normal, because the conjugation action of $S_d$ on $\KG$ by coordinate permutation is non-trivial — for $k \in \KG$ and $\sigma \in S_d$, the components of $k - \sigma\!\cdot\!k$ are $k_i - k_{\sigma^{-1}(i)} \in G$, which need not lie in $H$. Equivalently, $H^d \rtimes S_d$ is not normal in $(H^d + \KG) \rtimes S_d$. So the naive set-theoretic quotients $(\KG \rtimes S_d)/(K_H \rtimes S_d)$ and $((H^d + \KG) \rtimes S_d)/(H^d \rtimes S_d)$ are not groups, even though $\KG/K_H \cong K_{G/H}$ and $(H^d + \KG)/H^d \cong K_{G/H}$ as abstract groups.

% A precise description of the structure these morphisms carry would require concepts not introduced in this thesis, namely:
% \begin{itemize}
%     \item the \emph{free locus} $\cP^{[d], \circ} \subseteq \cP^{[d]}$, defined as the preimage of the complement of the big diagonal $\Delta \subseteq X^{(d)}$ (the open locus on which the residual $S_d$-action is free);
%     \item \emph{twisted forms} of a constant finite group, obtained as the contracted product $T \times^{S_d} A$ of an étale $S_d$-torsor $T$ with a finite group $A$ carrying an $S_d$-action; this is a finite étale group scheme, étale-locally isomorphic to the constant group $\underline{A}$ but not globally so when the $S_d$-action on $A$ is non-trivial and $T$ does not split;
%     \item \emph{torsors under non-constant group schemes}, generalising the constant-group torsors used throughout this thesis.
% \end{itemize}
% With these in hand, one shows that, restricted to the free locus, $q$ is a torsor under a twisted form of the constant group $K_{G/H}$, of degree $|K_{G/H}| = |G/H|^{d-1}$; the left vertical inherits the same description by base change. We do not need this refinement in what follows: only the four labelled torsors and the Cartesian property are used.
\end{rem*}

% \Cref{prop:torsor_tower_diagram} encodes two distinct factorizations of $\cP^{[d]_H} \to X^{(d)}$, which we shall refer to as the two \emph{towers}:

% \begin{description}
%     \item[The $H$-symmetric tower,] obtained by traversing the top edge and the right column of the diagram:
%     \[
%     \cP^{[d]_H} \xrightarrow{\;H\;} (\cP/H)^{(d)} \xrightarrow{\;q\;} (\cP/H)^{[d]} \xrightarrow{\;G/H\;} X^{(d)};
%     \]
%     \item[The $G$-symmetric tower,] obtained by traversing the left column and the bottom row of the diagram:
%     \[
%     \cP^{[d]_H} \longrightarrow \cP^{[d]_G} \xrightarrow{\;H\;} (\cP/H)^{[d]} \xrightarrow{\;G/H\;} X^{(d)}.
%     \]
% \end{description}
% In each tower the labelled arrows are torsors under the indicated constant groups; the unlabelled arrow and the morphism $q$ are not, in general, torsors under a constant group scheme.

% The Cartesian property in \Cref{prop:torsor_tower_diagram} expresses that the connection morphism $\cP^{[d]_H} \to \cP^{[d]_G}$ on the left is the base change of $q : (\cP/H)^{(d)} \to (\cP/H)^{[d]}$ along the structural map $\cP^{[d]_G} \to (\cP/H)^{[d]}$.
% Properties of $q$ that are stable under base change therefore descend automatically to the connection morphism on the left.

\section{Setup and notation}

Let $\cP \to U$ be a $G$-torsor over $U = C \setminus \mdls$, where $G$ is a
finite abelian group, $H \trianglelefteq G$, and $C$ is a smooth projective
curve. Fix $p_\infty \in \mdls$ and an integer $d$.
Assume the $G$ torsor $\cP \to U$ has ramification strictly bounded by $d$ at $p_\infty$. (i.e. $< d$)
according to the definition in \Cref{definition:local_ramification_boundness}
(i.e. the higher d'th-ramification group of the extension of DVRs at $p_\infty$ is trivial).

Recall the four quotient schemes, defined in:
\Cref{subsection:quotient_torsors_towers}
\[
\cP^{[d]_H},\qquad \cP^{[d]_G},\qquad (\cP/H)^{(d)},\qquad (\cP/H)^{[d]},
\]
together with the canonical Cartesian square

\[
\begin{tikzcd}[row sep=large, column sep=large]
\cP^{[d]_H} \ar[r,"H"] \ar[d]
  & (\cP/H)^{(d)} \ar[d,"q"] \\
\cP^{[d]_G} \ar[r,"H"'] \ar[dr,"G"',bend right=15]
  & (\cP/H)^{[d]} \ar[d,"G/H"] \\
  & U^{(d)},
\end{tikzcd}
\]
in which the arrows labelled $H$, $G/H$ and $G$ are torsors under the
indicated finite constant groups (in particular, \emph{\'etale}), and the
upper square is Cartesian (\Cref{prop:torsor_tower_diagram}).

The $G/H$-torsor $\cP/H \to U$ compactifies to a finite map of smooth
projective curves $f \colon C' \to C$. Pick $p'_\infty \in f^{-1}(p_\infty)$

\[
f^{(d)} \colon C'^{(d)} \to C^{(d)}
\]
for the induced map of $d$-fold symmetric powers; it sends
$d \cdot p'_\infty \mapsto d \cdot p_\infty$. Let
\[
\pi_C \colon X_C \to C^{(d)},\qquad \pi_{C'} \colon X_{C'} \to C'^{(d)}
\]
be the blowups at $d \cdot p_\infty$ and $d \cdot p'_\infty$ respectively, with
exceptional divisors $E_C,E_{C'}$ and generic points $\eta_C,\eta_{C'}$. The
DVRs of these generic points are
\[
V_C := \mathcal{O}_{X_C,\eta_C} \subset K(C^{(d)}),\qquad
V_{C'} := \mathcal{O}_{X_{C'},\eta_{C'}} \subset K(C'^{(d)}).
\]

\subsection*{A base-change lemma}

We will use the following standard lemma when comparing a Galois extension
with the function field of a base change.

\begin{lemma}\label{lem:galois-base-change}
Let $E/F$ be a finite Galois extension with group $H$, and let $K/F$ be an
arbitrary field extension. Fix a common overfield $\Omega \supseteq E,K$ and
form inside it the intersection $M := E \cap K$, the compositum $E \cdot K$,
and the subgroup $N := \Gal(E/M) \leq H$. Then:
\begin{enumerate}
  \item $E \cdot K/K$ is Galois, and restriction induces an isomorphism
        \[
        \Gal(E \cdot K/K) \xrightarrow{\ \sim\ } \Gal(E/M)=N;
        \]
  \item as $K$-algebras,
        \[
        E \otimes_F K \cong \prod_{i=1}^{[M:F]} E \cdot K.
        \]
\end{enumerate}
In particular, $E \otimes_F K$ is a field if and only if $M=F$, equivalently,
if and only if $E$ and $K$ are linearly disjoint over $F$.
\end{lemma}

\begin{proof}
Write $E=F[\alpha]$ and let $g\in F[x]$ be the minimal polynomial of $\alpha$.
Then $\deg(g)=|H|$, and $g$ splits completely in $E$ because $E/F$ is
Galois. The classical translation theorem shows that $E\cdot K/K$ is Galois
and that restriction identifies its Galois group with the subgroup of $H$
fixing $E\cap K=M$ pointwise, namely $\Gal(E/M)=N$.

For the second assertion, factor $g=g_1\cdots g_m$ into monic irreducible
polynomials over $K$. Adjoining any root of $g$ to $K$ produces $E\cdot K$,
so each $g_i$ has degree
\[
    [E\cdot K:K]=[E:M]=|N|.
\]
It follows that $m=|H|/|N|=[M:F]$, and hence
\[
E\otimes_F K
  \cong K[x]/(g)
  \cong \prod_{i=1}^{m} K[x]/(g_i)
  \cong \prod_{i=1}^{[M:F]} E\cdot K.
\]
This product is a field exactly when $m=1$, or equivalently when $M=F$.
\end{proof}

\section{The master diagram}\label{sec:master}

The torsor tower over the open base $U^{(d)}$ embeds into the projective
geometry of $C^{(d)}$ and $C'^{(d)}$, which in turn admit canonical blowups
at the points at infinity. We assemble all of this into a single commutative
diagram.

\[
\begin{tikzcd}[row sep=2.4em, column sep=2.2em]
\cP^{[d]_H} \ar[r,"H"] \ar[d]
  & (\cP/H)^{(d)} \ar[d,"q"] \ar[r,hookrightarrow,"\iota'"]
  & C'^{(d)} \ar[dd,"f^{(d)}"]
  & X_{C'} \ar[l,"\pi_{C'}"'] \ar[dd,dashed,"f_X"] \\
\cP^{[d]_G} \ar[r,"H"'] \ar[dr,"G"',bend right=10]
  & (\cP/H)^{[d]} \ar[d,"G/H"]
  & & \\
  & U^{(d)} \ar[r,hookrightarrow,"\iota"']
  & C^{(d)}
  & X_C \ar[l,"\pi_C"]
\end{tikzcd}
\]

Reading the diagram:
\begin{itemize}
  \item The \textbf{left block} (columns 1--2) is the torsor tower of
        \Cref{prop:torsor_tower_diagram}: $H$-torsors horizontally, the
        $G/H$-torsor on the lower right, and the composite $G$-torsor
        $\cP^{[d]_G} \to U^{(d)}$.
  \item The \textbf{middle block} (column 3) records the compactifications.
        The map $\iota$ is the canonical open immersion
        $U^{(d)} \hookrightarrow C^{(d)}$. The map $\iota'$ is the open
        immersion $(\cP/H)^{(d)} \hookrightarrow C'^{(d)}$ obtained from
        compactifying $\cP/H \subset C'$. The vertical $f^{(d)}$ is the
        finite morphism of symmetric powers; restricting along $\iota'$
        recovers the canonical map $(\cP/H)^{(d)} \to U^{(d)}$, which is
        the composite of $q$ with the $G/H$-torsor.
  \item The \textbf{right block} (column 4) records the blowups at the
        diagonal points at infinity. The dashed arrow
        $f_X \colon X_{C'} \dashrightarrow X_C$ is the rational lift of
        $f^{(d)}$, defined on the strict transform; on generic points of
        the exceptional divisors it sends $\eta_{C'} \mapsto \eta_C$ so we have
        induced extension of DVRs $V_C \hookrightarrow V_{C'}$         
\end{itemize}

\section{Unramifiedness of $\cP^{[d]_G}$ at $\eta_C$}
So, returning to the beginning of this appendix, a natural attempt to prove
part~(3) of \Cref{theorem:cyclic_cases_after_blowup} from part~(2) by
induction on $r$ comes down
to proving the following theorem. 

%% [REVIEW: unlabeled -- this is the central theorem of the appendix, referred to repeatedly in prose ('the Theorem'); add a label and use \Cref]
\begin{theorem}
Suppose
\begin{enumerate}
  \item[\textnormal{(H1)}] $(\cP/H)^{[d]} \to U^{(d)}$ is unramified at
        $\eta_C$; and
  \item[\textnormal{(H2)}] $\cP^{[d]_H} \to (\cP/H)^{(d)}$ is unramified
        at $\eta_{C'}$.
  \item [\textnormal{(H3)}]$G=\Z/p^r\Z$ is cyclic of prime power order.
\end{enumerate}
Then $\cP^{[d]_G} \to U^{(d)}$ is unramified at $\eta_C$.
\end{theorem}

And the next section is devoted to explaining why this is not an easy task, 
and in fact, we were not able to prove it.

\section{Difficulties in Proving the Theorem}

First, note that we can assume $\cP \to U$ is connected, and that $H$ is non-trivial. 
\begin{enumerate}
\item[\textnormal{(R1)}] \emph{$\cP \to U$ is connected.}  Let
      $\rho \colon \pi_1(U) \to G$ be the monodromy of $\cP$ and put
      $G_0 := \mathrm{Im}(\rho)$, $H_0 := H \cap G_0$.  Pick a connected
      component $\cP_0 \subseteq \cP$; it is a connected $G_0$-torsor, and
      $\cP$ is the disjoint union of $[G : G_0]$ copies of $\cP_0$ permuted
      by $G$.  Since the symmetric-power quotients of
      \Cref{subsection:quotient_torsors_towers} commute with such disjoint
      unions, each of the four schemes in \Cref{prop:torsor_tower_diagram}
      for $(\cP, G, H)$ decomposes as a disjoint union of copies of the
      corresponding scheme for $(\cP_0, G_0, H_0)$.  Unramifiedness at
      $\eta_C$ (resp.\ $\eta_{C'}$) is detected componentwise, so (H1) and
      (H2) are inherited by $(\cP_0, G_0, H_0)$; (H3) is inherited because
      subgroups of cyclic $p$-groups are cyclic $p$-groups.  Replacing
      $(\cP, G, H)$ by $(\cP_0, G_0, H_0)$ we may assume $\cP$ is connected
      and $\rho$ is surjective.
\item[\textnormal{(R2)}] \emph{$H \neq 0$.}  If $H = 0$ then
      $\cP/H = \cP$ and $(\cP/H)^{[d]} = \cP^{[d]_G}$, so (H1) is
      literally the conclusion.
\end{enumerate}

We can read a decomposition of fields directly from the canonical Cartesian square of
\Cref{prop:torsor_tower_diagram}. After the reduction (R1) all four corners are integral:
the three corners $\cP^{[d]_G}$, $(\cP/H)^{[d]}$ and $(\cP/H)^{(d)}$ are integral,
and the fourth, $\cP^{[d]_H} = (K_H \rtimes S_d)\backslash \cP^{\times_S d}$, is integral as
a quotient of the integral scheme $\cP^{\times_S d}$
(\Cref{theorem:int_geoint_implies_productint}). So we may take function fields; using the open
immersions $\iota,\iota'$ of the master diagram together with
$K_C = \Frac V_C = K(C^{(d)}) = K(U^{(d)})$ and
$K_{C'} = \Frac V_{C'} = K(C'^{(d)}) = K\bigl((\cP/H)^{(d)}\bigr)$, we identify
\[
K(U^{(d)}) = K_C,\qquad
K\bigl((\cP/H)^{[d]}\bigr) = F,\qquad
K\bigl((\cP/H)^{(d)}\bigr) = K_{C'},\qquad
K\bigl(\cP^{[d]_G}\bigr) = E,
\]
with $E/F$ Galois of group $H$, $F/K_C$ Galois of group $G/H$, and the composite $E/K_C$
Galois of group $G$. The fourth corner $\cP^{[d]_H}$ is the fibre product
$\cP^{[d]_G} \times_{(\cP/H)^{[d]}} (\cP/H)^{(d)}$; being integral, its single function
field is the compositum $E\cdot K_{C'}$. In particular its generic fibre $E \otimes_F K_{C'}$
over $(\cP/H)^{(d)}$ is a field, so by \Cref{lem:galois-base-change} the extensions $E$ and
$K_{C'}$ are linearly disjoint over $F$; that is,
\[
E \cap K_{C'} = F ,
\]
and $E\cdot K_{C'}/K_{C'}$ is Galois with group $H$. We denote:
\begin{enumerate}
    \item $K_C := \Frac V_C$, $K_{C'} := \Frac V_{C'}$;
    \item $F := K\bigl((\cP/H)^{[d]}\bigr)$, $E := K(\cP^{[d]_G})$;
    \item $V_F := V_{C'}|_F$;
    \item a place $v_E$ of $E$ above $V_F$, and a place $v_\star$ of $E \cdot K_{C'}$
          above both $v_E$ and $V_{C'}$.
\end{enumerate}

Thus we have:
\[
\begin{tikzcd}[row sep=1.6em, column sep=1.4em]
 & E \cdot K_{C'} \arrow[dl, dash] \arrow[dr, dash, "H"] & \\
E \arrow[dr, dash, "H"'] & & K_{C'} \arrow[dl, dash] \\
 & F \arrow[d, dash, "G/H"] & \\
 & K_C &
\end{tikzcd}
\qquad\qquad
\begin{tikzcd}[row sep=1.6em, column sep=1.4em]
 & v_\star \arrow[dl, dash] \arrow[dr, dash] & \\
v_E \arrow[dr, dash] & & V_{C'} \arrow[dl, dash] \\
 & V_F \arrow[d, dash] & \\
 & V_C &
\end{tikzcd}
\]
(lines denote field inclusion / place lying-over, read upwards; in the left
diagram a label is the Galois group of the extension where it is Galois).  In
addition $E/K_C$ is Galois with group $G$.

Read through this dictionary, the hypotheses of the theorem become statements about
the ramification of the places above:
\begin{itemize}
  \item[\textnormal{(H1)}] $(\cP/H)^{[d]} \to U^{(d)}$ is unramified at $\eta_C$
        means that the $G/H$-extension $F/K_C$ is unramified at $V_C$. In particular
        $V_F$ is unramified over $V_C$, with $e(V_F \mid V_C) = 1$.
  \item[\textnormal{(H2)}] $\cP^{[d]_H} \to (\cP/H)^{(d)}$ is unramified at $\eta_{C'}$
        means that the $H$-extension $E\cdot K_{C'}/K_{C'}$ is unramified at $V_{C'}$;
        equivalently the inertia group $I(v_\star \mid V_{C'}) \leq H$ is trivial.
  \item[\textnormal{(H3)}] $G = \Z/p^r\Z$ makes $G$, hence $H$ and every inertia group,
        a cyclic $p$-group; as the residue characteristic is $p$, the ramification of
        $E/F$ is potentially \emph{wild}.
\end{itemize}

The conclusion we are after --- that $\cP^{[d]_G} \to U^{(d)}$ is unramified at $\eta_C$ ---
is the statement that the $G$-extension $E/K_C$ is unramified at $V_C$, i.e.\ that the inertia
$I(v_E \mid V_C) \leq G$ is trivial. By (H1) the lower layer $F/K_C$ is already unramified at
$V_C$, so this inertia lies in $H = \Gal(E/F)$ and coincides with the inertia of the top layer,
\[
I(v_E \mid V_C) \;=\; I(v_E \mid V_F) \;\leq\; H .
\]
The theorem therefore reduces to the single assertion
\[
E/F \text{ is unramified at } V_F, \qquad\text{i.e.}\qquad I(v_E \mid V_F) = 1 .
\]

Now $E\cdot K_{C'}/K_{C'}$ is the base change of $E/F$ along $F \hookrightarrow K_{C'}$, and under
the restriction isomorphism $\Gal(E\cdot K_{C'}/K_{C'}) \cong \Gal(E/F) = H$ the inertia injects
\[
I(v_\star \mid V_{C'}) \;\hookrightarrow\; I(v_E \mid V_F) .
\]
Hypothesis (H2) says the source is trivial; the goal is that the target is trivial. These need
\emph{not} coincide: the two differ precisely by the ramification of $K_{C'}/F$ at $V_{C'} \mid V_F$
--- the ramification that the symmetric-power base change $q$ introduces at the point at infinity
--- which can be wild and can absorb the ramification of $E/F$. It is exactly this discrepancy
that obstructs deducing the goal from (H1)--(H3).

\textcolor{red}{Put in other words:
The base change $K_{C'} /F$ is wildly ramified, so the inertia comparison is not onto .}

\subsection*{The inertia injection, and why \textnormal{(H2)} is not enough}

The injection used above is an instance of the following general fact.

\begin{lemma}\label{lem:inertia-injection}
Let $E/F$ be a finite Galois extension, linearly disjoint over $F$ from a field extension
$K/F$, with compositum $E\cdot K$; thus restriction is an isomorphism
$r \colon \Gal(E\cdot K/K) \xrightarrow{\ \sim\ } \Gal(E/F)$, $\sigma \mapsto \sigma|_E$.
Let $w$ be a place of $E\cdot K$, and write $w_K := w|_K$, $w_E := w|_E$, $w_F := w|_F$ for its
restrictions. Then $r$ carries the inertia group of $w$ over $w_K$ into that of $w_E$ over $w_F$:
\[
  r\bigl(I(w \mid w_K)\bigr) \;\subseteq\; I(w_E \mid w_F),
\]
so that $I(w\mid w_K) \hookrightarrow I(w_E\mid w_F)$ is injective. Moreover, completing at $w$
identifies $I(w\mid w_K)$ with the inertia group of the local compositum $\widehat E\cdot\widehat K$
over $\widehat K$; in particular, if $\widehat E \subseteq \widehat K$ then $I(w\mid w_K) = 1$,
no matter how ramified $E/F$ is at $w_F$.
\end{lemma}

\begin{proof}
That $r$ is an isomorphism is the translation theorem for the Galois extension $E/F$ and the
linearly disjoint $K/F$ (\Cref{lem:galois-base-change}(1) with $M = F$). For the inclusion, let
$\sigma \in I(w\mid w_K)$. Since $E/F$ is Galois and $\sigma$ fixes $F \subseteq K$ pointwise,
$\sigma(E) = E$ and $\sigma|_E \in \Gal(E/F)$ fixes $w_F$. As $\sigma$ lies in the decomposition
group of $w$ we have $w\circ\sigma = w$; restricting to $E$ gives $w_E\circ\sigma|_E = w_E$, so
$\sigma|_E \in D(w_E\mid w_F)$. Finally the residue extension
$\kappa(w_E)\hookrightarrow\kappa(w)$ is equivariant and $\sigma$ acts trivially on $\kappa(w)$,
hence $\sigma|_E$ acts trivially on $\kappa(w_E)$, i.e.\ $\sigma|_E\in I(w_E\mid w_F)$.
Injectivity is inherited from that of $r$. For the last assertion, the completion of $E\cdot K$
at $w$ is the local compositum $\widehat E\widehat K$, with $\Gal(\widehat E\widehat K/\widehat K)
= D(w\mid w_K)$ and inertia subgroup $I(w\mid w_K)$; if $\widehat E\subseteq\widehat K$ then
$\widehat E\widehat K = \widehat K$ and this group is trivial.
\end{proof}

Taking $w = v_\star$ (so $w_K = V_{C'}$, $w_E = v_E$, $w_F = V_F$) recovers the injection
$I(v_\star\mid V_{C'}) \hookrightarrow I(v_E\mid V_F)$. The implication ``(H2) $\Rightarrow$ goal''
would follow if this were always onto; the next example shows it need not be, already for
$G = \Z/p^2\Z$ over $\F_p$.

\begin{example}\label{ex:p2-counterexample}
Take $k = \F_p$, $G = \Z/p^2\Z$, and $H = p\,\Z/p^2\Z \cong \Z/p\Z$ the unique subgroup of order
$p$, so $G/H \cong \Z/p\Z$. We describe the local fields at the place $V_{C'}\mid V_F$ --- the only
place where the obstruction lives --- writing $\widehat{(\,\cdot\,)}$ for completions.

\smallskip
\noindent\emph{The base and the $G$-tower.} Let $\widehat{K_C} = \F_p((t))$, and present the local
$G$-extension $\widehat E/\widehat{K_C}$ (an Artin--Schreier--Witt extension) by its two layers:
\begin{itemize}
  \item lower layer $\widehat F/\widehat{K_C}$ (group $G/H$): \emph{split}, so $\widehat F = \F_p((t))$.
        Then $F/K_C$ is unramified at $V_C$ --- this is \textnormal{(H1)} --- and $\widehat{V_F} = \widehat{V_C}$;
  \item upper layer $\widehat E/\widehat F$ (group $H$): the Artin--Schreier extension
        \[
          \widehat E = \widehat F(y), \qquad y^p - y = t^{-1}.
        \]
        By \Cref{theorem:artin_schreier_ramification}(3) (with $m = 1$, prime to $p$) this is totally,
        wildly ramified of degree $p$, with $I(v_E \mid V_F) = H \cong \Z/p\Z \neq 1$.
\end{itemize}
Hence $I(v_E\mid V_C) = I(v_E\mid V_F) = H \neq 1$: the extension $E/K_C$ is ramified at $V_C$, so
\emph{the conclusion of the theorem fails}.

\smallskip
\noindent\emph{The base change.} The places of $K_{C'}$ over $V_F$ are governed by the
symmetric-power cover $q$. Suppose at $V_{C'}$ this cover reproduces the \emph{same}
Artin--Schreier extension,
\[
  \widehat{K_{C'}} = \widehat F(y),\quad y^p - y = t^{-1}, \qquad\text{i.e.}\qquad \widehat{K_{C'}} = \widehat E .
\]
This is compatible with the global linear disjointness $E\cap K_{C'} = F$: two \emph{distinct}
degree-$p$ extensions of the global field $F$ may have isomorphic completions at a place. Now the
local compositum is $\widehat{E\cdot K_{C'}} = \widehat E\cdot\widehat{K_{C'}} = \widehat E
= \widehat{K_{C'}}$, so $E\cdot K_{C'}/K_{C'}$ is split, hence unramified, at $V_{C'}$:
\[
  I(v_\star\mid V_{C'}) = \Gal\bigl(\widehat E/\widehat{K_{C'}}\bigr) = 1 .
\]
Thus \textnormal{(H2)} holds, and \textnormal{(H3)} holds by construction.

\smallskip
\noindent\emph{Outcome.} All of \textnormal{(H1)--(H3)} hold, yet $E/K_C$ is wildly ramified at
$V_C$ with inertia $H\cong\Z/p\Z$. The inertia injection of \Cref{lem:inertia-injection} is here
\[
  \underbrace{I(v_\star\mid V_{C'})}_{=\,1}\;\hookrightarrow\;\underbrace{I(v_E\mid V_F)}_{=\,\Z/p\Z},
\]
the inclusion of the trivial group: (H2) carries no information about the ramification of $E/F$,
which has been entirely \emph{absorbed} by that of $K_{C'}/F$ --- the wild ramification introduced
by $q$ at the point at infinity.
\end{example}

\begin{rem*}
This is globally consistent with everything above. Since $\widehat E\otimes_{\widehat F}\widehat{K_{C'}}
= \widehat E\otimes_{\widehat F}\widehat E$ is not a field, $\cP^{[d]_H}$ (function field
$E\cdot K_{C'}$) \emph{splits} at $V_{C'}$ into several places and is therefore unramified at
$\eta_{C'}$, as \textnormal{(H2)} demands; meanwhile $\cP^{[d]_G}$ stays ramified at $\eta_C$. The
hypotheses simply do not see the local coincidence $\widehat{K_{C'}} = \widehat E$, and ruling out
such coincidences for a given $(\cP, d)$ is exactly the global input the inductive argument lacks
--- which is why this is a wrong path.
\end{rem*}

\subsection*{Absorption made explicit: two local fields}

\Cref{ex:p2-counterexample} posited the local coincidence $\widehat{K_{C'}} = \widehat E$ abstractly.
For completeness we exhibit it in the smallest possible setting --- two honest local fields and a
one-line Artin--Schreier reduction --- to make the absorption phenomenon fully visible.

\begin{example}[absorption of a wild break, explicit fields]\label{ex:absorption-explicit}
Let $k = \bar{\F}_p$ and let $F = k((t))$. Define a second local field $K = k((\pi))$, and embed
$F \hookrightarrow K$ by
\[
t \;=\; \frac{\pi^{p}}{1-\pi^{p-1}} \;=\; \pi^{p}\bigl(1+\pi^{p-1}+\pi^{2(p-1)}+\cdots\bigr),
\qquad\text{so}\qquad v_\pi(t)=p .
\]
Since $t = \pi^p\cdot(\text{unit})$ and $\,dt/d\pi = -\pi^{2p-2}(1-\pi^{p-1})^{-2}\neq 0$, the extension
$K/F$ is separable, \emph{totally (wildly) ramified of degree $p$}; equivalently it is the
Artin--Schreier extension $\sigma^p-\sigma = t^{-1}$ with $\sigma=\pi^{-1}$, of break $1$
(\Cref{theorem:artin_schreier_ramification}(3)).

Now let $E/F$ be the Artin--Schreier extension
\[
E = F(y),\qquad y^p - y = t^{-1}.
\]
Over $F$ this is ramified: $v_t(t^{-1}) = -1$ is prime to $p$, so by
\Cref{theorem:artin_schreier_ramification}(3) it is totally wildly ramified of degree $p$, break $1$,
with $I(v_E\mid v_F) = \Z/p\Z \neq 1$.

\smallskip
\noindent\emph{Yet $E$ is absorbed by $K$.} Computing $t^{-1}$ inside $K$,
\[
t^{-1} \;=\; \pi^{-p}\bigl(1-\pi^{p-1}\bigr) \;=\; \pi^{-p}-\pi^{-1} \;=\; \wp\bigl(\pi^{-1}\bigr),
\qquad \wp(x):=x^p-x ,
\]
so $y=\pi^{-1}\in K$ already solves $y^p-y=t^{-1}$. Hence $E = F(\pi^{-1}) \subseteq K$ (indeed
$E = K$, both being degree $p$ over $F$), the compositum is $E\cdot K = K$, and
\[
E\cdot K / K \ \text{is trivial, hence unramified:}\qquad I(v_\star\mid v_K) = 1 ,
\]
while $E/F$ is ramified. The ramification of $E/F$ has been \emph{entirely absorbed} by the wild
base change to $K$ --- exactly the failure of surjectivity in the inertia injection
$I(v_\star\mid v_K)\hookrightarrow I(v_E\mid v_F)$, here $1 \hookrightarrow \Z/p\Z$.
\end{example}

\begin{rem*}
\emph{The mechanism, and why $F$ cannot do it.} Read over $K$, the datum $t^{-1}$ has pole order
$v_\pi(t^{-1}) = p$, \emph{divisible by $p$}; this is what licenses the Artin--Schreier reduction
$y\mapsto y-\pi^{-1}$, which removes the pole completely because $\wp(\pi^{-1})$ matches $t^{-1}$ to
all orders. The substitution is available over $K$ but \emph{not} over $F$: the reducing element
$\pi^{-1}$ has valuation $-1$, whereas $v_\pi(F^\times) = p\,\Z$ contains no element of valuation
$-1$ at all. Thus the wild base change does not ``fight'' the ramification of $E/F$; it
\emph{supplies the one element of intermediate valuation} that the Artin--Schreier reduction
needs. The ramification index $e(K/F)=p$ is the gate that turns the prime-to-$p$ break $1$ into the
divisible-by-$p$ pole order $p$ over $K$.

\emph{General $r$.} For $G=\Z/p^r$ the faithful local model has $\widehat{K_{C'}}/\widehat F$ totally
ramified of degree $p^{r-1}$ (the residue field $k$ being algebraically closed, there is no residual
part), and absorption is the proper inclusion $\widehat E \subsetneq \widehat{K_{C'}}$ of the
top-layer $\Z/p$-extension. The explicit fields are obtained by iterating the recipe above: with
$E = k((\pi))$, $\pi=y^{-1}$ as here, set $\pi = \rho^{p}(1-\rho^{p-1})^{-1}$ to get a degree-$p$
overfield $k((\rho)) \supsetneq E$, and so on $r-1$ times. The pole order of the top datum is
multiplied to $p^{r-1}\cdot 1 \equiv 0 \pmod p$ and reduced away by the same one-line substitution.
Over an algebraically closed residue field, ``$E\cdot K/K$ unramified with $E/F$ ramified'' is
\emph{equivalent} to $\widehat E\subseteq\widehat{K_{C'}}$ --- there are no nontrivial unramified
extensions to hide in --- so absorption is precisely a containment of local fields, never a softer
phenomenon.
\end{rem*}
